\section{From integral schedules to tensor intertwiners}
\label{sec:columns}

The switch integrality in Definition~\ref{def:schedule} is weaker than
the usual integral saturated-chain condition in a highest-weight path
model.  We now prove the representation-theoretic implication that is
needed here directly, using fundamental columns.

\begin{proposition}[Column intertwiners]\label{prop:columns}
Let $\lambda,\mu,\xi$ be dominant integral weights of type $D_n$,
$n\ge2$.  Choose a shape obtained by concatenating straight segments
whose increments are fundamental weights and whose sum is $\mu$.
If this shape admits an integral schedule from $\lambda$ to $W\xi$,
then
\[
 \Hom_G\bigl(V(\lambda)\otimes V(\mu),V(\xi)\bigr)\ne0.
\]
An empty list of columns is allowed when $\mu=0$.
\end{proposition}

We first record carefully what happens to integrality when a schedule
is folded into the dominant chamber.

\begin{lemma}\label{col:folding}
An integral schedule can be folded into a dominant path with Weyl
direction states in decreasing Bruhat order.  Its changes of state
can be expressed as an ordered sequence of downward reflections,
including at coincident times.  At each such step, from $w$ to
$s_\alpha w$ at time $t$,
\begin{equation}\label{col:folded-integrality}
 \alpha\cdot w\eta(t)\in\Z.
\end{equation}
For a concatenation of fundamental columns, all column endpoint
positions of the folded path are dominant integral weights.
\end{lemma}

\begin{proof}
Use the simplifications and folding construction in the proof of
Lemma~\ref{lem:schedule-moment}.  After deleting removable switches,
perturb the remaining strict schedule to separate its sites and make
ordinary chamber crossings transverse.  In the perturbed path, an
original switch becomes a downward reflection in the current chamber
frame, and a sorting crossing becomes a downward simple reflection.
There are at most the longest Weyl length many such strict downward
steps.  Taking a subsequence fixes their discrete labels and order;
their times then converge.  This gives the asserted ordered chain even
if several steps acquire the same limiting time.  The perturbations
used for this construction need not be integral.

Integrality is checked on the limiting chain, in its specified order.
For an original switch, changing the chamber frame changes both its
root and its direction by the same Weyl element.  Its pairing in
\eqref{col:folded-integrality} is therefore the original integral
pairing.  For a sorting crossing, suppose all preceding coefficients
are integral and write the position before that crossing as
\begin{equation}\label{col:prefix}
 x=\lambda+w\eta-\sum_{j< i}m_j\beta_j.
\end{equation}
The wall equation $\alpha\cdot x=0$ gives
\[
 \alpha\cdot w\eta
 =-\alpha\cdot\lambda
   +\sum_{j<i}m_j\,\alpha\cdot\beta_j\in\Z.
\]
The usual switch identity preserves \eqref{col:prefix} after the
crossing.  Induction therefore proves integrality for every prefix,
also inside a limiting cluster.  At a column endpoint $\eta$ is a
weight, so \eqref{col:prefix} shows that $x$ is a weight.  It is
dominant by construction.
\end{proof}

Consider a single column of weight $\omega$.  Normalize its parameter
to $0\le\tau\le1$, so that the original shape on this interval is
$\chi+\tau\omega$, with $\chi$ a sum of preceding fundamental
weights.  Write $x_0,x_1$ for its dominant endpoints and put
$\delta=x_1-x_0$.  Let $u$ and $v$ be the first and last Weyl states
in its open interior.  Boundary-time steps are left between the
columns; thus the last state of one column dominates the first state
of the next.  Since $\alpha\cdot w\chi$ is integral,
Lemma~\ref{col:folding} gives at every interior step
\begin{equation}\label{col:local-level}
 \tau\,\alpha\cdot w\omega\in\Z.
\end{equation}

\paragraph{Minuscule columns.}
For a minuscule weight, all root pairings are $0$ or $\pm1$.
Consequently \eqref{col:local-level} permits no interior reflection
that changes $w\omega$.  The column is straight, with
$\delta=u\omega$.  The extremal case of
Lemma~\ref{lem:coefficient} gives an intertwiner
\[
 V(x_0)\otimes V(\omega)\longrightarrow V(x_1)
\]
whose highest-to-highest coefficient is nonzero at the fixed flag $u$.
This includes the vector column and both spin columns.  For $D_2$
and $D_3$ all fundamental columns are minuscule.
The classical exterior and half-spin models are described in
\cite[Lectures~19--20]{FH91}; here only their root-pairing property
is needed for the minuscule case.
Although full Weyl states inside the column may change while fixing
$\omega$, this causes no problem for composition: the chosen fixed
waypoint has future $u$, and $u\ge v$ dominates the next column's
initial state.

It remains to treat an exterior column
\[
 \omega=e_1+\cdots+e_k,\qquad 2\le k\le n-2.
\]
Let $E$ be the standard $2n$-dimensional quadratic space, with basis
$E_{+i},E_{-i}$ satisfying
$(E_{+i},E_{-j})=\delta_{ij}$ and all other basis pairings zero.
The module $\bigwedge^k E$ is the fundamental module $V(\omega)$;
its flag vectors are wedges of isotropic $k$-planes.

\paragraph{The coefficient to be constructed.}
For the exterior column, Lemma~\ref{lem:coefficient} reduces the
intertwiner problem to a functional supported at weight $\delta$.
The following criterion isolates its annihilation conditions. We will
then construct such a functional and a flag at which it is nonzero,
with past and future limits $u,v$ so that columns can be composed.

\begin{lemma}\label{col:common-walls}
Let $J$ be the set of simple roots $\alpha_i$ with
$x_0\cdot\alpha_i=x_1\cdot\alpha_i=0$.
A functional supported on the weight space of weight $\delta$ in
$\bigwedge^k E$, invariant under the semisimple group generated by
$J$, is an allowed highest-to-highest coefficient for an intertwiner
$V(x_0)\otimes V(\omega)\to V(x_1)$.
\end{lemma}

\begin{proof}
Put $h_i=x_0\cdot\alpha_i$ and $h'_i=x_1\cdot\alpha_i$.
By Lemma~\ref{lem:coefficient}, the functional must annihilate
$f_i^{h_i+1}V(\omega)$ for every $i$.  A vector contributing to its
weight-$\delta$ component would have source weight
$\delta+(h_i+1)\alpha_i$, whose simple-root pairing is
\[
 \bigl(\delta+(h_i+1)\alpha_i\bigr)\cdot\alpha_i
 =h_i+h'_i+2.
\]
All weights of $\bigwedge^k E$ have coordinates in $\{0,1,-1\}$,
so this pairing cannot exceed two.  Since $h_i,h'_i$ are nonnegative
integers, the source is absent unless both vanish.  In that case
$\alpha_i\in J$, and invariance annihilates $f_iV(\omega)$.
\end{proof}

\paragraph{The midpoint and a subgroup waypoint.}
The nonzero root pairings of $\omega$ have absolute value at most two.
By \eqref{col:local-level}, an effective interior switch occurs only
at $\tau=\tfrac12$ and has absolute pairing two.  It flips the signs
on two occupied axes.  A root with pairing zero preserves both the
occupied support and its complement.  Thus the support $S$ of
$w\omega$ is fixed throughout the column, every interior reflection
belongs to $D_S\times D_{S^c}$, and
\begin{equation}\label{col:endpoint-directions}
 a=u\omega,\qquad a'=v\omega,\qquad
 \delta=\frac{a+a'}2.
\end{equation}
In particular, the set
\[
 R=\{i\in S:a'_i=-a_i\}
\]
has even cardinality.  On $S\setminus R$ one has
$\delta_i=a_i=a'_i\in\{1,-1\}$, whereas $\delta_i=0$ on $R$.

Let $H$ be the connected subgroup for $D_S\times D_{S^c}$.
The orbit $H\cdot uB$ is a flag variety of $H$: its stabilizer at
$uB$ is the Borel $H\cap uBu^{-1}$, and the ambient positive Borel
induces the positive Borel $H\cap B$.  Each interior comparison is a
downward Bruhat reflection in this orbit.  Indeed its root is positive
for $H\cap B$, and its inverse image by the current state is negative
for the corresponding isotropy Borel, exactly as in the ambient
comparison.  The resulting Bruhat inequality supplies a waypoint
\begin{equation}\label{col:waypoint}
 b\in U_H^+u\cap U_H^-v
   \subset U^+u\cap U^-v.
\end{equation}
The two unipotent groups here are the positive and negative ones
induced from the ambient group.  Equivalently, one may obtain this
intersection by the Bruhat bridge construction used in
Lemma~\ref{lem:flag-composition}.

The wedge at $b$ represents a maximal isotropic plane in
$E_S=\operatorname{span}\{E_{+i},E_{-i}:i\in S\}$.
Its $a$ and $a'$ coordinates are nonzero, since their lines are its
past and future limits.  Use the $a$ coordinate as a graph chart:
write its basis as
\[
 z_j=E_{a_j j}+\sum_{i\in S}T_{ij}E_{-a_i i}\quad(j\in S).
\]
Here the notation $E_{a_i i}$ denotes $E_{+i}$ or $E_{-i}$ according
to the sign $a_i$.  Isotropy says that $T$ is skew-symmetric.  The
$a'$ coordinate of $\bigwedge_{j\in S}z_j$ is, up to sign,
$\det T_{R,R}$, so
\begin{equation}\label{col:invertible-principal}
 T_{R,R}\text{ is nonsingular}.
\end{equation}
When $R$ is empty, the empty determinant is understood to be one.

There is a useful restriction on its entries:
\begin{equation}\label{col:earlier-negative}
 i<j,\quad T_{ij}\ne0\quad\Longrightarrow\quad a_i=-1.
\end{equation}
Indeed a single replacement in the wedge gives a nonzero coordinate
of weight $a-a_i e_i-a_j e_j$.  Since $a$ is the past extremal
limit in \eqref{col:waypoint}, the difference
$-a_i e_i-a_j e_j$ has positive pairing with the regular dominant
flow direction.  This difference is a root, so it is positive.
In type $D$, a root supported on
$i<j$ is positive precisely when its $i$-coefficient is positive,
which proves \eqref{col:earlier-negative}.

\paragraph{A common-wall invariant functional.}
We now have a waypoint with the required flow limits and a
nonsingular principal matrix $T_{R,R}$.  We first construct a
weight-$\delta$ functional satisfying Lemma~\ref{col:common-walls}.
The free coefficients in this construction will then be chosen to
make its value at the waypoint nonzero.

The sign change on the last coordinate preserves the positive root
system, interchanging its two end simple roots.  Apply it, if needed,
to arrange $x_0\ge0$ coordinatewise; this changes none of the
preceding arguments.  Let $\mathcal B$ be the collection of blocks
of axes having $\delta_i=0$ and the same positive value of $x_{0i}$.
Assign a scalar $c_B$ to each such block, put $c_i=c_B$ for $i\in B$,
and put $c_i=0$ on the zero-position, zero-$\delta$ tail.  Define
\begin{equation}\label{col:theta}
 \Theta=\sum_{B\in\mathcal B}c_B
       \sum_{i\in B}E_{+i}^*\wedge E_{-i}^*,
 \qquad
 \Psi=\left(\bigwedge_{i:\,\delta_i\ne0}
                  E_{\delta_i i}^*\right)
              \wedge\Theta^{|R|/2}.
\end{equation}
Choose any fixed order for the coordinate wedge.  The functional
$\Psi$ has degree $k$ and is supported on the weight $\delta$.

It is invariant under the common-wall semisimple group.  To see this
directly, a common wall $e_i-e_{i+1}$ requires equal $x_0$ coordinates
and equal $\delta$ coordinates.  On a block with nonzero $\delta$,
the coordinate wedge in \eqref{col:theta} is the determinant
covector and is invariant under the special-linear group.  On a block
with positive $x_0$ and zero $\delta$, the two-form
$\sum E_{+i}^*\wedge E_{-i}^*$ is the canonical invariant pairing
form for the standard space and its dual, hence is invariant under
the whole general-linear group.  The remaining zero-position,
zero-$\delta$ tail carries a special-orthogonal factor and contributes
no covectors.  If the other end root $e_{n-1}+e_n$ belongs to $J$, then
$x_{0,n-1}=x_{0,n}=0$.  Dominance of $x_1$ gives
$(\delta_{n-1},\delta_n)=(0,0)$ or $(1,-1)$.  The first case belongs
to the orthogonal zero tail.  In the second, the difference end root
is absent, and changing the last sign gives a special-linear
determinant block.  These cases cover the group generated by $J$.

\paragraph{Nonzero evaluation at the waypoint.}
It remains to choose the $c_B$ so that $\Psi$ is nonzero at $b$.
The coordinate factors in \eqref{col:theta} select exactly their own
primal columns in the graph chart.  After these factors are removed,
the pullback of $\Theta$ to the remaining columns indexed by $R$ has
skew matrix
\begin{equation}\label{col:modified-skew}
 \widetilde T_{ij}=(a_i c_i+a_j c_j)T_{ij},\qquad i,j\in R,
\end{equation}
up to one overall sign convention.  If $\zeta_i$ are the coordinate
functions on the primal plane, then
$E_{+i}^*\wedge E_{-i}^*$ pulls back to
$a_i\,\dd\zeta_i\wedge\sum_jT_{ij}\dd\zeta_j$;
summing the two contributions to an entry gives
\eqref{col:modified-skew}.  The top exterior power is nonzero
precisely when $\widetilde T$ is nonsingular.  This is automatic for
$R=\varnothing$; assume henceforth that $R\ne\varnothing$.

Partition $R=N\sqcup P$, where
\[
 N=\{i\in R:x_{0i}>0,\ a_i=-1\}.
\]
Thus $P$ consists of positive signs at positive positions, together
with all indices in the zero tail.  Dominance of the first part of
the column gives dominance of $x_0+\eps a$ for all sufficiently
small $\eps>0$.  Within an equal positive-position block, the signs
therefore occur in decreasing order.  In the zero tail no negative
sign can precede another index: only the final coordinate may be
negative.  Combined with \eqref{col:earlier-negative}, this proves
\begin{equation}\label{col:skew-blocks}
 T_{R,R}=\begin{pmatrix}A&B_0\\-B_0^{\mathsf T}&0\end{pmatrix}
 \quad\text{in the ordering }N,P.
\end{equation}
Furthermore a nonzero entry from $i\in N$ to a positive-position
$j\in P$ must have $i<j$ and $x_{0i}>x_{0j}$.  Equality of the
positions would contradict the decreasing sign order in that block.
For a zero-tail $j\in P$, every $i\in N$ already precedes $j$,
regardless of the sign of a possible final coordinate.

Order the blocks in $\mathcal B$ from larger to smaller position.
Choose positive exponents $s_B$ strictly decreasing in this order,
write $B_i$ for the block containing a positive-position index $i\in R$,
and set
\[
 c_B=1+t^{s_B}\qquad(t>0).
\]
We let $t\to0^+$.  For a permitted $N,P$ entry with both positions
positive, the larger-position block has the larger exponent, so its
perturbation decays faster.  This separates the blockwise coefficients
while leaving their common limit equal to one.
In \eqref{col:modified-skew}, the $N,N$ block tends to $-2A$.
For a positive-position $P$ index $j\in B$, scale its row and column
by $t^{-s_B}$.  On a permitted entry from $i\in N$ to this $j$,
the multiplier becomes
\[
 t^{-s_B}(-c_i+c_j)
 =1-t^{s_{B_i}-s_B}\longrightarrow1.
\]
For a zero-tail $P$ index leave the row and column unscaled and then
negate both: its unscaled multiplier is $-c_i\to-1$.
The $P,P$ block stays zero.  Thus the scaled modified matrix tends to
\begin{equation}\label{col:skew-limit}
 \begin{pmatrix}-2A&B_0\\-B_0^{\mathsf T}&0\end{pmatrix}.
\end{equation}
This matrix is nonsingular.  More generally, if
$\gamma\ne0$ and $(x,y)$ lies in the kernel after replacing $A$ in
\eqref{col:skew-blocks} by $\gamma A$, then $(x,y/\gamma)$ lies in
the kernel of the original matrix.  Nonsingularity therefore follows
from \eqref{col:invertible-principal}.  By continuity,
$\widetilde T$ is nonsingular for sufficiently small positive $t$.
This proves that the functional $\Psi$ evaluates nontrivially on the
waypoint wedge.  Lemma~\ref{col:common-walls} now supplies the
required single-column intertwiner.

\begin{proof}[Proof of Proposition~\ref{prop:columns}]
Apply Lemma~\ref{col:folding} and perform the preceding construction
on each column.  Write its endpoints as $h_{j-1},h_j$, its weight
as $\omega_j$, and choose the resulting intertwiner
\[
 V(h_{j-1})\otimes V(\omega_j)\longrightarrow V(h_j).
\]
Its highest-to-highest coefficient is nonzero on the chosen waypoint.
For an exterior column the past and future flags are its first and
last states $u_j,v_j$; for a minuscule column both may be taken to
be $u_j$.  At a column boundary, $v_j\ge u_{j+1}$.
Thus the chosen future flag dominates the next past flag: directly
for an exterior column, and through
$u_j\ge v_j\ge u_{j+1}$ for a minuscule column.
Lemma~\ref{lem:flag-composition} therefore makes
the composition nonzero on the Cartan summand
\[
 V\Bigl(\sum_j\omega_j\Bigr)=V(\mu)
 \subset\bigotimes_jV(\omega_j).
\]
The first and last dominant positions are $h_0=\lambda$ and
$h_q=\xi$, proving the assertion.  For an empty list the folded
path is constant, so $\lambda=\xi$ and the identity map suffices.
\end{proof}
